% !PRINTED 5
\weberpage{5}
The task that will now occupy us first is a twofold one: by the introduction of new variables the differential
\weq{5}{\frac{dx}{\sqrt{f(x)}} \quad \text{or} \quad \frac{y\,dx-x\,dy}{\sqrt{f(x,y)}}}
is to be transformed into another of the same form, but one in which the function under the square-root sign is made as simple as possible and, in particular, made to depend on as small a number of parameters as possible (§ 3, 4, 5). Secondly, the general differential (3) or (4) is to be decomposed into other similar ones in which the function $\Phi(x)$ or $\Phi(x,y)$ is as simple as possible (§ 12).

The function $f(x)$ can be decomposed into linear factors (Bd. I, § 33). We accordingly set also
\[
f(x) = a_0(x-x_1)(x-x_2)(x-x_3)(x-x_4).
\]
Let us think of the quantities $x, x_1, x_2, \ldots$ as abscissas of points laid off on a straight line from an arbitrary origin, the positive ones towards the one side, the negative ones towards the other (for example, positive to the right, negative to the left). Then the differences $x-x_1, x-x_2, x-x_3, x-x_4$ represent the distances of the point $x$ from $x_1, x_2, x_3, x_4$, and indeed reckoned positive when $x$ lies towards the positive side of $x_1, x_2, \ldots$. In the same sense $x_2-x_1$ is the distance of the point $x_2$ from $x_1$, and so on.

We also denote the differences more briefly by $(xx_1), (xx_2), \ldots (x_2x_1), \ldots$ We use this mode of representation only for the sake of an easier survey, and we do not exclude imaginary points of the line.

§ 2. \gesperrt{Cross Ratios.}

The simplest means for the simplification of the elliptic differential rests upon the linear transformation of the function $f(x)$ (Bd. I, § 67). Let $\alpha, \beta, \gamma, \delta$ denote four arbitrary constants whose determinant
\weq{1}{\alpha\delta - \beta\gamma = r}
is different from zero. We set
\weq{2}{x = \frac{\alpha x' + \beta}{\gamma x' + \delta},}
%CONT>>
